\documentclass[17pt, aspectratio=169]{beamer}

% --- Paquetes de Idioma y Estética ---
\usepackage[utf8]{inputenc}
\usepackage[T1]{fontenc}
\usepackage[spanish]{babel}
\usepackage{xcolor}
\usepackage{tcolorbox}
\tcbuselibrary{skins}

% --- Definición de Colores ICMG ---
\definecolor{ICMGBlue}{HTML}{002366}
\definecolor{ICMGGold}{HTML}{996515}
\definecolor{CitationBg}{HTML}{F4F4F4}

% --- Configuración del Tema Beamer ---
\setbeamercolor{palette primary}{bg=ICMGBlue, fg=white}
\setbeamercolor{palette secondary}{bg=ICMGBlue, fg=white}
\setbeamercolor{structure}{fg=ICMGBlue}
\setbeamercolor{title}{fg=ICMGBlue}
\setbeamerfont{title}{series=\bfseries}

\setbeamertemplate{navigation symbols}{}

\setbeamertemplate{footline}{
  \leavevmode%
  \hbox{%
  \begin{beamercolorbox}[wd=.5\paperwidth,ht=2.25ex,dp=1ex,left]{author in head/foot}%
    \hspace*{2ex} Clase 2: El Arrepentimiento
  \end{beamercolorbox}%
  \begin{beamercolorbox}[wd=.5\paperwidth,ht=2.25ex,dp=1ex,right]{date in head/foot}%
    Iglesia Cristiana Manantial de la Gracia \hspace*{2ex} \insertframenumber{} / \inserttotalframenumber \hspace*{2ex} 
  \end{beamercolorbox}}%
  \vspace{2pt}
}

% --- Formato para Citas Bíblicas ---
\newtcolorbox{citaBiblicaSlide}[1]{
  enhanced,
  colback=CitationBg,
  colframe=ICMGGold,
  leftrule=3pt,
  rightrule=0pt,
  toprule=0pt,
  bottomrule=0pt,
  arc=0mm,
  fonttitle=\bfseries\color{ICMGBlue},
  title={#1},
  boxrule=0pt,
  italic
}

\title{Clase 2: El Arrepentimiento}
%\subtitle{Clase 1: El Fundamento de Nuestra Identidad}
\author{Iglesia Cristiana Manantial de la Gracia}
\date{Serie: El Bautismo}

% --- Configuración de Diapositiva de Título Intermedia ---
\AtBeginSection[]{
  \begin{frame}[plain] % 'plain' quita el pie de página para resaltar el título
    \centering
    \vfill
    {\color{ICMGGold}\rule{0.6\textwidth}{1.5pt}} % Línea decorativa superior
    \vspace{0.4cm}
    
    {\color{ICMGBlue}\scshape\large SECCIÓN \thesection \par} % "SECCIÓN 1", etc.
    \vspace{0.2cm}
    {\color{ICMGBlue}\Huge\bfseries \insertsectionhead \par} % Nombre de la sección
    
    \vspace{0.4cm}
    {\color{ICMGGold}\rule{0.6\textwidth}{1.5pt}} % Línea decorativa inferior
    \vfill
  \end{frame}
}

\begin{document}

\begin{frame}[plain]
    \titlepage
    \centering
    \color{ICMGGold}\rule{0.5\textwidth}{1pt}
\end{frame}

\section{¿Qué es el arrepentimiento?}

\begin{frame}{¿Qué es el arrepentimiento?}
    \begin{itemize}
        \item Cambio radical de mentalidad y propósito.
    \end{itemize}
\end{frame}

\begin{frame}{¿Qué es el arrepentimiento?}
    \begin{itemize}
        \item Cambio radical de mentalidad y propósito.
        \item Involucra:
        \begin{itemize}
            \item Verdadero sentido de culpa y pecaminosidad (2 Cor. 7:10).
            \item Abrazar la misericordia de Dios en Cristo.
            \item Aborrecimiento por el pecado y giro a Dios (Salmo 119:128).
            \item Anhelo constante de una vida santa y obediente a Dios.
        \end{itemize}
    \end{itemize}
\end{frame}

\section{¿Qué NO es el arrepentimiento?}

\begin{frame}{¿Qué NO es el arrepentimiento?}
    \begin{itemize}
        \item Miedo a las consecuencias (Hechos 8:18-24).
    \end{itemize}
\end{frame}

\begin{frame}{¿Qué NO es el arrepentimiento?}
    \begin{itemize}
        \item Miedo a las consecuencias (Hechos 8:18-24).
        \item Reconocer el pecado (1 Samuel 15:17-21, 24-25).
    \end{itemize}
\end{frame}

\begin{frame}{¿Qué NO es el arrepentimiento?}
    \begin{itemize}
        \item Miedo a las consecuencias (Hechos 8:18-24).
        \item Reconocer el pecado (1 Samuel 15:17-21, 24-25).
    \end{itemize}

    \begin{citaBiblicaSlide}{Salmo 51:17}
        "Los sacrificios de Dios son el espíritu contrito;
Al corazón contrito y humillado, oh Dios, no despreciarás."
    \end{citaBiblicaSlide}
    
\end{frame}

\begin{frame}{}

    \begin{citaBiblicaSlide}{Marcos 1:15}
        "«El tiempo se ha cumplido», decía, «y el reino de Dios se ha acercado; arrepiéntanse y crean en el evangelio»."
    \end{citaBiblicaSlide}
    
\end{frame}

%\begin{frame}{¿Qué es el bautismo?}
    %\begin{itemize}
        %\item Medio de Gracia (no salva).
        %\item Reconocer que la Biblia es la única %autoridad y estándar infalible.
    %\end{itemize}
    
    %\begin{citaBiblicaSlide}{Juan 17:3}
        %"Y esta es la vida eterna: que te conozcan a ti, %el único Dios verdadero, y a Jesucristo, a quien %has enviado."
    %\end{citaBiblicaSlide}
    
%\end{frame}


\section{El Hijo Pr\'odigo}
%PAsajes y cómo cumple con la definición, Lucas 15:11-24

\begin{frame}{El Hijo Pr\'odigo: Lucas 15:11-24}
    \begin{itemize}
        \item Reconoce su miseria (v. 17).
        
    \end{itemize}
\end{frame}

\begin{frame}{El Hijo Pr\'odigo: Lucas 15:11-24}
    \begin{itemize}
        \item Reconoce su miseria (v. 17).
        \item Abandono de su condici\'on para ir con su padre (v. 20).
        
    \end{itemize}
\end{frame}

\begin{frame}{El Hijo Pr\'odigo: Lucas 15:11-24}
    \begin{itemize}
        \item Reconoce su miseria (v. 17).
        \item Abandono de su condici\'on para ir con su padre (v. 20).
        \item Reconoce ofensa y pecado (v. 21).

    \end{itemize}
\end{frame}

\begin{frame}{Conclusi\'on}
    \begin{itemize}
        \item El fruto del arrepentimiento lo vemos en nuestro pecados m\'as atesorados.
    \end{itemize}
\end{frame}

\begin{frame}{Conclusi\'on}
    \begin{itemize}
        \item El fruto del arrepentimiento lo vemos en nuestro pecados m\'as atesorados.
        \item Agradezcamos por el pecado del que Dios nos ha librado.
    \end{itemize}
\end{frame}

\begin{frame}{Conclusi\'on}
    \begin{itemize}
        \item El fruto del arrepentimiento lo vemos en nuestro pecados m\'as atesorados.
        \item Agradezcamos por el pecado del que Dios nos ha librado.
        \item Dependamos de Dios para vivir en obediencia.
    \end{itemize}
\end{frame}

\end{document}